% -----------------------------------------------------------
\section{Talk Ideas}
% -----------------------------------------------------------
	\label{TALKIDEAS}
	
	\begin{description}
	
		\item[View Lehi/Nephi/promised land story and celestial kingdom]
			View this stody as a type of our trying to enter into the highest degree of the celestial kingdom (1N)
			Requires being married (1N 7:1); requires great sacrifice; knowledge from the Lord; death and resurrection (into the water/coming out of it), etc.
			See 1N 8:1, seeds; other verses that mention seed(s).
			Connect to Lehi's vision, the fruit desirable to make one happy? And connect that to 2N 5:27 and 2N 9:43? Lots of other connections to make here, in his dream and the way he describes the fruit and everything.
				So you see Lehi's dream as an allegorical expression of what is happening to his family as they journey to the new land; and then you see both of those as allegorical expressions of our journey to the highest degree of the celestial kingdom.
				Lehi wants his FAMILY to partake of it.
				Water in both the vision and the journey to the new land (the sea).
				Tree in the vision; they take seeds for ``fruit of every kind'' (8:1).
				Lehi doesn't go get them personally after taking the fruit; he beckons and calls to them only. Doesn't approach L/L either; he apparently just beckons and calls but they refuse.
				The role of the SEEDS they carry.
				
		\item[Track each testimony gained in 1 Nephi]
			the gaining of each person's testimony, and the differences in experiences that each person has as they gain their testimonies (1N)
			Nephi and Sam, early. Sariah, when her sons return safely. Laman and Lemuel, perhaps earlier, but they say ``we know of a surety'' after Nephi shocks them in 1N 17:52--55, especially 55. What about Zoram? And the members of Ishmael's family? Track each person, very interesting.
			Note that Sariah uses the term ``surety'' in 5:8. Remember Lehi struggles with the food---each person has their own struggles and their own moments of conviction based on what is important to them.
			
		\item[The rod in Lehi's dream]
			In Lehi's dream, how we go through the rod/temptation/fruit cycle over and over, and how in different aspects of our life we are on different places in different paths holding on to different ``rods''
			
		\item[The tree of life in Lehi's dream with other trees]
			Specifically with the tree of life as mentioned at the end of Alma 32, and how the word grows into that tree. Note also the connection to John 4 and the well of water springing up. I think this gives you a very new view on the tree in Lehi's vision and the possibilities for interpreting it. Remember with this, the word in Lehi's vision is the rod, but it's the seed in Alma. Interesting to compare.
			
		\item[Nephi's discussion of ``plainness'' at various places]
			He uses the word a lot and is really into it. You could search for ``plain'' and ``plainness'' and see how he uses it.
			
		\item[The differences between how Limhi's and Alma's people escape]
			I'm specifically interested in \hyperref[MOSIAH-23-23]{Mosiah 23:23} and \hyperref[MOSIAH-24-21]{24:21}, which explicitly says that no one except the Lord could deliver them; I'll have to check again, but I don't think this is said about Limhi's people. So what's the difference, and what differences does it cause in what the people do? Why are some situations such that the Lord is the only one who can deliver? How can we tell what situation we're in, if at all, and how should it change our behavior? See also \hyperref[MOSIAH-29-19]{Mosiah 29:19--20}.
			
		\item[Compare and contrast the journeys across the water (1N/Ether)]
			Note that Nephi's group goes crazy (18:9) and faces lots of setbacks in the weather (18:13); the boj's group praised the Lord the entire time (Eth 6:9; compare Nephi's attitude in 1N 18:16) but still faced lots of challenges (Eth 6:6--7). The main difference seems to be that the Lord caused the wind to blow the entire time in the boj's favor (Eth 6:5, 8), while with Nephi's group, because of their wickedness I'm guessing, they actually went in reverse, away from their goal (18:13). Note the contrast between 18:15, they're about to die, and Eth 6:10, no monster could break them, and they had light the entire time.
			Another good parallel: the boj's people had light the entire time, but for Nephi's group, their source of light was the Liahona, and it stopped working (18:12) until Nephi took charge again (18:21). Since they didn't have the light of the compass, they didn't know where to steer the ship, so they couldn't have even gotten there on their own if they had tried to fight the wind pushing against them.
			1N 17:3 kind of describes this situation: when we're righteous, the Lord prepares the way before us.
			So a few morals:
				\begin{compactitem}
				\item Nephi had the right attitude, but because of the people around him, he wasn't progressing.
				\item Being righteous isn't a guarantee of calm seas. But it is a guarantee that (1) we'll have the light of the Lord with us, and (2) the wind will always be pushing in the direction of the promised land, whatever happens.
				\item Sometimes, because of the decisions of other people, our prayers are totally ineffectual (1N 18:19).
				\end{compactitem}
	
		\item[{\hyperref[MOSIAH-2-9-NOTE-view]{``Views'' of God's mysteries in Benjamin's story}}]
			There is an interesting progression in what ``views'' of God's mysteries are unfolded to the people in the story of King Benjamin. Think of how the ``unfolding'' process occurs, where some parts that are visible are then no longer visible once a larger portion has unfolded; maybe some of those parts that aren't visible any longer, once we are ``viewing'' the good things in the future, are our sins and carnal nature.
			
		\item[{\hyperref[SIN-definition]{The definition of sin}}]
			Analyzing James 4:17 and what we learn from it about sin.
			
		\item[{\hyperref[SIN-sin-transgression]{The difference between sin and transgression}}]
			Embracing the growth mindset, which is like playing the game where sin is involved, versus looking for a commandment mindset, in which what you care about is fulfilling the commandments you have been given.
			
		\item[The relationship between joy and fun, or having fun]
			Thinking specifically of Pook's essay on the fountain of youth, and \hyperref[MOSIAH-3-19]{Mosiah 3:19}. Becoming like a child---not necessarily with the qualities listed in that verse, but with the much more salient qualities that children tend to have. This topic is related to the new birth, and gaining that childlike energy and zest for life is an important part of the new birth (and one of its primary goals!). See \hyperref[NEWBIRTH-child]{here} for more. See also \hyperref[JOHN-4-14]{John 4:14} and notes there. See also the below idea ``Becoming like a child in Mosiah 3:16--21.''
		
		\item[Becoming like a child in Mosiah 3:16--21]
			Focusing on a few things in particular:
				(1) the qualities of children mentioned in verses 16--18 and 21 as a context for verse 19---see \hyperref[MOSIAH-3-18-NOTE-children]{this note} to Mosiah 3:18;
				(2) the unusual list of qualities associated with being a child in verse 19, which are not the ones that we normally associate with children;
				(3) the fact that people damn themselves (verse 18) unless they become like little children. Why is this? One thing little children do is \textit{grow}, not just physically but mentally, spiritually, and emotionally, so becoming like a child involves adopting a growth mindset and working to reach an abundance mentality;
				
		\item[Limhi's people and Alma's people]
			Compare/contrast the hardships and escapes of Limhi's people (Mo 20--22) and those of Alma's people (Mo 23--24)
			How are they similar? How are they different? Why does the Lord act differently in each case? How do the people respond differently? Note Mo 25:10 as well, which talks about the power of God in delivering them only when mentioning Alma; doesn't mention Limhi's people and the power of God. But then see Mo 25:16.
		
		\item[15 key BOM scriptures (for an institute class sometime?]
			If I'm ever asked, I can imagine teaching an institute class where we focus on a single BOM verse every week, for a total of between 12--15 verse throughout the semester. Here are some ideas of verses such a course could include:
			
			\begin{enumerate}
				\item \hyperref[1NEPHI-4-6]{1 Nephi 4:6} (leading into a discussion of this entire chapter, I think; see footnotes in \hyperref[1NEPHI-4]{1 Nephi 4}.)
				\item \hyperref[1NEPHI-10-17]{1 Nephi 10:17} (maybe in conjunction with \hyperref[MORMON-9-21]{Mormon 9:21}?)
				\item \hyperref[MOSIAH-3-16]{Mosiah 3:16}
				\item \hyperref[MOSIAH-3-19]{Mosiah 3:19}
			\end{enumerate}
			
		\item[Retaining a remission of sin]
			Look at how this works in the scriptures, with a focus on \hyperref[MOSIAH-4-11]{Mosiah 4:11--12}.
			
		\item[Disruptions to our lives in 1 Nephi 16]
			Using the events of 1 Nephi 16, talk about how to respond to disruptions to our lives. Talk about how we must be in the discomfort zone in order to progress; talk about Nephi's ways of reacting to the disruption, and the ways of everyone else. Make a point that it is better to think about this story in terms of each character being a separate part of \textit{ourselves}---that is, each of us has inside of us a Nephi, a Lehi, a Laman, and Sariah, the daughters of Ishmael, etc. And we allow different characters to control our thoughts and actions and habits at different times. I also have some footnotes that are relevant here.
			
		\item[How do we apply deliberate practice to developing righteousness?]
		
		\item[The experience of Nephi's family as an example of developing self-mastery]
			Notice that they are in the wilderness for eight \textit{years}---and then they get to somewhere amazing! And then they still have to cross the waters because there is something \textit{even better}!
			
		\item[John needing assurance that Christ was the one they sought]
			See \hyperref[LCHRIST-088-su]{this note}.
			
		\item[The conversion of Lamoni as an almost literal rebirth and resurrection (Al 18--19)]
			This really starts at the end of 18, where Lamoni falls as if he were dead. He is raised as if from the dead by his wife---what significance does that have? Abish tells everyone the news; he gets up and ``stands'' (common when talking about the resurrection); many beautiful parallels here, including hand-to-hand contact as in the temple and through the veil; see also Al 47:22--24. A beautiful story with a beautiful parallel; this could be powerful when fleshed out. This is almost a literal rebirth for them. Note the word ``ground'' for the grave, parallels like that. The words are very particular. Note what the people say in 19:33, that they have no more desire to do evil. That's what Benjamin's people say in Mo 5:2. Those people didn't fall asleep or go unconscious, but they saw themselves ``in their own carnal state, even less than the dust of the earth.'' There are two cases of death language there---the dust, referring to the grave, and their carnal state, for being carnally minded is death (Rom 8:6, 2N 9:39). So I shouldn't only be looking for the rebirth process in the scriptures, I should also be looking at how that process is a type of the resurrection, and how the language used to describe one is often used to describe the other.
			See also what happens with Lamoni's father, especially in 22:18, 22.
			For more about being dead in sin and raised from it, see Eph 2:1--6.
			This also answers a question about Al 13:3, what a preparatory redemption is. It's the new birth. It's a \textit{preparatory} redemption because it isn't the big one that will happen when we are resurrected and glorified; but it's still a redemption, a passing from death unto life through a new birth. It's preparatory to the real one, in the same way that temple ordinances are preparatory to actual becoming a king/queen and priest/priestess. The new birth itself is a preparatory ``ordinance,'' so to speak (and maybe an actual one), that will be fulfilled and actualized when we are resurrected and sanctified.
			See also Al 22:13, where the plan of redemption gets us out of the carnal state.
			See also:
				DC 84:26--27 (preparatory gospel is the gospel of repentance and baptism, and remission of sins); 
				
		\item[Alma's song in Alma 29]
			Compare to the other ``songs'' in the scriptures, like Nephi and the others in the OT. Compare also to Nephi's song in Hel 7:7--9.
			
		\item[Compare what Korihor taught to what the Zoramites believed]
			Why would they have run him through and killed him? Compare what Alma et al. taught in response (Al 30--31; 32--34)	For the first part, see 30:13 and 31:22, for example. For the second part, compare what Alma teaches about the effects of the word being ``real'' because they are ``light'' (32:35), to what Korihor preaches in 30:16.
			
		\item[Discuss how one has to LEARN to live a spiritual life]
			As told by Alma to Helaman in Al 37:32--37. It may not come naturally---someone has to teach us how to do this. Related to deliberate practice and the gospel.
			
		\item[One way the Nephites fall]
			Slowly setting aside what they know until they are left to their own strength (Hel 4) I read this chapter today, 11/30/2015, and felt a definite rebuke. Lately I haven't been very spiritually diligent, and I think this chapter explains why. I just let things slide and got too casual, relied on myself too much. A good chapter to use to call people to repentance.
			This is also related to the Independent Study talk on being spiritually awake.
			
		\item[The brother of Jared delays four years in calling on the Lord (Eth 2)]
			Why did he wait four years, after they arrived at the water's edge? I bet he knew what was going to be required of him, and it was too much for him to think about; they were too happy there; his kids were growing up; they had already left behind other places that would have been good to live in; they get to the sea, and he's got to realize what the next step is. But he's scared and comfortable so he doesn't do anything about it.
			What's the next step for us spiritually? What sea are we facing down right now? What are we delaying or avoiding because we already know what's required of us, and how difficult it will be? The Lord didn't let them stop in other good places (Eth 2:7) because he had something better for them; what does the Lord have which is better for us, that we aren't receiving because we aren't calling on him to ask for his help in overcoming the last great hurdle?
			
		\item[The building of the barges to cross the sea (Eth 2)]
			The brother of Jared asks how they are going to \textit{steer} (19) without light. But they're inside the boat, and it's enclosed, so how could they steer whether they have light inside or not? By ``light'' he must mean, how are we going to know where we're going? This makes the story more interesting and the parallel to us more powerful---sometimes we're going to have to board barges in faith without being able to see where we're going at all. No light illuminates the forward path; all we can do is believe that we are being guided by the Lord in the proper way.
			Think of this story in comparison to Nephi's story of crossing the waters---the entire thing, from preparing the boat (Nephi had to learn how to build it/brother of Jared used his old boats as a model), to steering (Nephi had to steer vs. enclosed ship; see also 1N 18:13 vs. Eth 2:19), and so on.
			Notice the difference in the questions he asks about light---in 19 it's about steering, while in 22 it's about having light inside the vessel so they don't have to cross in darkness.
			Follow up on this with the information that it's 6:1--10.
			
		\item[The way the people cross the waters in Ether 6]
			They are tossed all over chaotically, but are overall on a course toward some place better.
			
		\item[Use the preparing of the stones as a metaphor to describe how the Lord answers questions]
			Or how he gives a testimony, or how human writings become divine scripture---he takes what is already there and makes it shine with inspiration and divine light and intelligence (Eth 2--3)
			In our lives we are sometimes swallowed up in the depths of the sea (2:25). But the scriptures are a source of light for us, prepared to illuminate us when nothing else will. How do human writings become inspired? Humans take of the material that they already have, like stones, and think and refine and try to understand things as best they can. Then they approach the Lord with their questions and thoughts, each a precious stone that they have worked on refining. The Lord touches each one, causing each one to bring forth light and have power to guide on its own----he corrects mistakes; he sends a testimony; he causes ideas to shine with power that they didn't have before, by adding intelligence through the spirit; and so on. We take what we have to the Lord and he makes it shine.
			
		\item[The Lord appears to the brother of Jared (Eth 3)]
			A few points to note:
			\begin{compactitem}
				\item In 6 it doesn't say that the Lord appeared or arrived; he was already there, talking to the brother of Jared in person. But because of the veil he couldn't see him until a certain moment. Extremely powerful---how close is God to us when we pray? Are there times that he's right next to us, and we just can't discern him? He just stretched forth his hand, already being right there, and the the boj's perception changed. 
				\item The pre-earth Christ, not having received a body, is able to physically interact with the stones. Not immaterial.
				\item The veil is taken off his \textit{eyes}, not his mind or anything else. There's something about the perceptual faculty itself, beginning with the eye, that needs to be changed in order for us to perceive God. The physical matter and processes have to be changed, quickened maybe, in order for us to perceive the Lord. Note verses 19--20 here as well; it may be that the boj had gained so much knowledge, had so much experience with the Lord, that the veil had already been effectively removed from his mind, and what was blocking him know was the physical matter and processes of his perceptual capacities. The mental or intelligent component was already removed from the veil because he had already worked with the Lord so much.
				\item Note also that right before this the boj had been chastised by the Lord for not calling on him for four years (!). Then soon after that chastising experience, the Lord asks him how he wants to cross the waters with light (showing confidence in him), and then revealing himself to him (manifesting his worthiness). Extremely important I think.
			\end{compactitem} 
			
		\item[The ``order of the son of God'' (DC 107)]
			An interesting way to begin a discussion about the priesthood (in this section and elsewhere) could be to ask, why is the real name of the priesthood ``the holy priesthood, after the order of the Son of God,'' and not ``the holy priesthood, after the order of God''? See my notes there for expansions on this idea, including \hyperref[JOHN-5-26]{John 5:26}, the stuff about the \hyperref[LAW-christ]{law of Christ}, and so on. All the \hyperref[PRIESTHOOD]{priesthood} stuff.
			
		\item[The metaphor of frost, dew, and gathering (DC 121)]
			See the notes at \hyperref[DC121-11]{verse 11} and \hyperref[DC121-45]{verse 45}; this parallel is worth looking at closely, both in the section and in the letter. Think carefully about what dew is, how it gathers, how you get a sizable amount of water from it, and so on.
			
		\item[The ways of knowing, according to JS]
			Compare JS's ways of coming to know God and spiritual things vs. those of, for example, Ibn Rushd in the \textit{Decisive Treatise}; see \hyperref[DC46-14]{DC 46:14} and notes there.
			
		\item[My experience with what I thought I did to Keaton]
			Sin and memory.
			
		\item[Nephi being ``constrained'' to say certain things by the Spirit]
			It happens various times, and is worth studying; also worth studying how little other prophets are ``constrained'' in the same way. I think it tells you something about how Nephi experienced the Holy Ghost.
			
		\item[Nephi and Jared families crossing the ship in boats built in unusual ways]
			The Lord reveals to them how to do it, and they do it in ways that are not familiar to those around them.
			
		\item[\gr{οἱ δὲ εὐθέως ἀφέντες} in Matthew 4:20, 22]
			Talk about this phrase and how it's the same in both verses. If these verses were written about us, what would come next? What would we be leaving? Note that the last parts of the verses are identical too.
			
		\item[``If thou wilt''---``I will'' in Matthew 8:2--3]
		
		\item[How are we bound? (Matthew 12:29 and the later verses in Matthew 12]
			The later ones are about the man leaving his house after it's been swept out and cleaned. See also \hyperref[MARK-3-27]{Mark 3:27}.
			
		\item[Approaching Christ without the intent to make something happen, without faith]
			See notes at the priesthood document for \hyperref[MARK-5-24]{Mark 5:24--34}.
			
		\item[Failed preaching/missions leading to greater things; failure required first]
			See footnotes at \hyperref[HELAMAN-13-2]{Helaman 13:2}.
			
		\item[How the relationship between Jared and his brother changes over time as the brother develops spiritually]
	
		\item[Differences and similarities between how Nephi and the brother of Jared solve their beach and boat problems]
			See my footnote at \hyperref[ETHER-2-14]{Ether 2:14}.
			
		\item[On our contribution to the revelatory process in DC 9:11]
			See my notes there.
	
	\end{description}